Rust~\cite{matsakis2014rust} is designed to be a \emph{type-sound} language,
aiming to prevent various undesirable behaviors, including memory bugs like
use-after-free, at compile time through type checking~\cite{jung2017rustbelt}.
%
Recognizing this benefit, researchers and practitioners have adopted Rust for
systems programming to enhance software reliability.
%
Operating systems, web browsers, and network stacks have been developed in
Rust~\cite{levy2017multiprogramming, lankes2020rustyhermit,
narayanan2020redleaf, boos2020theseus, hu2024unishyper, peng2024framekernel,
dai2024verifying, hong2024taming, chiang2024securing, anderson2016engineering,
chen2025atmosphere, shang2025lwrustip};
%
legacy systems, including Linux, have begun integrating Rust~\cite{li2024rust,
netstack3, rust-dropbox, rust-coreutils};
%
and techniques to translate legacy code into Rust have been extensively
studied~\cite{emre2021translating, emre2023aliasing, zhang2023ownership,
hong2023concrat, hong2024dont, hong2024tag, wu2025genc2rust, hong2025forcrat,
hong2025automatically, hong2024type, zhou2025c2rusttv, xu2025optimizing,
zhang2025systematic, cai2026rustmap, zhang2025scalable}.
%
All these lines of work rely on the fundamental assumption that the Rust type
checker is sound.

Unfortunately, \rustc{}, the official Rust compiler, is not free from defects.
%
Liu et al.~\cite{liu2025empiricala} recently demonstrated that the compiler
suffers from various bugs, including crashes, miscompilations, and improper
diagnostic messages.
%
While all bugs warrant attention, \emph{soundness bugs}, where the type checker
accepts a program that should be rejected, are especially critical, as they
undermine Rust's core value of type soundness guarantee.
%
Notably, among the 301 \rustc{} bugs collected by Liu et al., 22 were identified
as soundness bugs.

Despite their importance, soundness bugs in the \rustc{} implementation have not
yet been studied in depth.
%
For instance, it remains unclear which language features are most prone to
soundness bugs, even though such insights would allow researchers to prioritize
specific features when developing testing techniques or formalizing the type
system.
%
While Liu et al. analyzed \rustc{} bugs and provided valuable insights, their
study encompassed all bug categories rather than focusing specifically on
soundness bugs.
%
We believe that an in-depth study dedicated to soundness bugs can benefit the
community and complement existing work.

In this work, we study 30 \emph{soundness issues} collected from the \rustc{}
GitHub issue tracker, each representing a case where a developer reported that
the compiler accepted code that should not pass type checking.
%
We analyze each issue in depth to characterize the following:
%
the \emph{affected feature},
%
the \emph{symptom} (i.e., how the feature is erroneously handled),
%
the \emph{consequence} (i.e., the specific undesirable behavior enabled by the
issue),
%
the features that \emph{trigger} the issue,
%
the \emph{consensus} of the community regarding whether it is indeed a bug,
%
and the \emph{lifecycle} of the issue, encompassing its introduction, discovery,
and fix.

Furthermore, we investigate artifacts that can aid in testing the type soundness
of \rustc{}.
%
Liu et al.'s study showed that existing testing tools for
\rustc{}~\cite{fuzz-rustc, tree-splicer, icemaker, sharma2023rustsmith,
yang2024rusttwins, wang2024rustlantis, dewey2015clp} revealed crash bugs but
failed to identify a single soundness bug.
%
A primary reason for this is the absence of a \emph{test oracle} capable of
determining whether a given program should be accepted or rejected.
%
To address this, RustSmith~\cite{sharma2023rustsmith} employs differential
testing, comparing results across different \rustc{} versions and
configurations.
%
However, its effectiveness is limited because most soundness issues persist
across multiple versions and settings, as we observe in our study.
%
This underscores the necessity of a test oracle tailored to type soundness.

For this reason, we study the following implementations, which serve as
potential test oracles:
%
AddressSanitizer~\cite{serebryany2012addresssanitizer} (a memory bug detector),
%
Miri~\cite{jung2026miri, gh-miri} (a Rust interpreter for detecting undefined
behavior),
%
Chalk~\cite{gh-chalk} (an implementation of the Rust trait system based on a
logic solver),
%
and a-mir-formality~\cite{gh-amirformality} (an executable model of the Rust
type system).
%
We assess their effectiveness by checking whether they can determine that the
programs from the collected issues should be rejected.

In addition, we study existing documentation that explains the semantics of
Rust.
%
Previous research has demonstrated that \emph{mechanized
specifications}~\cite{park2021jiset, youn2024bringing} can serve as effective
test oracles for language implementations,
%
detecting numerous bugs in JavaScript engines and
transpilers~\cite{park2021jest, park2023feature, ryu2024javascript} and
WebAssembly engines~\cite{youn2025west}, as well as soundness bugs in the P4
type checker~\cite{lee2026failing}.
%
Although Rust currently lacks a full specification~\cite{are-we-spec-yet},
documents such as the Rust Reference~\cite{rust-reference}, the FLS~\cite{fls},
and Rust RFCs~\cite{rust-rfcs} provide explanations of the semantics.
%
Therefore, we investigate the current state of this documentation and assess
whether these explanations are sufficiently clear to develop test oracles for
type soundness.

% ~\cite{rust-reference},
% ~\cite{rust-rfcs}
% ~\cite{rustc-dev-guide}

\vspace{0.5em} {\bf Contributions.} Overall, our contributions are as follows:
%
\begin{itemize}[leftmargin=*]
  %
  \item We construct a three-and-a-half-year dataset of 30 soundness issues in
    \rustc{} (\cref{sec:method}).
  %
  \item We conduct an empirical study of the soundness issues across multiple
    dimensions, including their affected features, symptoms, consequences,
    triggers, consensus, and lifecycles (\cref{sec:bugs}).
  %
  \item Based on our analysis, we summarize key findings and provide actionable
    suggestions for researchers and \rustc{} developers to better detect and
    mitigate soundness bugs (\cref{sec:findings}).
  %
  \item We assess existing artifacts that could potentially serve as oracles for
    testing type soundness, discussing their current capabilities and providing
    guidelines for their use (\cref{sec:oracles}).
  %
\end{itemize}

{\bf Summary of Findings.} We summarize our key findings as follows:
%
\begin{itemize}[leftmargin=*]
  %
  \item Certain soundness bugs lead to memory bugs; specifically, \emph{implied
    bounds} and \emph{trait objects} are critical triggering features that can
    compromise memory safety.
  %
  \item \emph{Associated types} frequently serve as a mechanism for exploiting
    edge cases in other features, and the interaction between \emph{lifetimes}
    and \emph{traits} poses challenges for sound type checking.
  %
  \item Many bugs have persisted since the initial introduction of the relevant
    features, typically requiring a long time to discover and remaining present
    across multiple compiler versions.
  %
  \item AddressSanitizer and Miri can detect soundness bugs that manifest as
    memory bugs; other bug categories could be identified by a-mir-formality and
    Chalk in the future, but they are currently immature.
  %
  \item Existing documents are often unclear, and constructing a mechanized
    specification for Rust cannot be achieved simply by collecting them.
  %
\end{itemize}
